% -----------------------------------------------------------
\section{Idle}
% -----------------------------------------------------------
	\label{IDLENESS}
	
% % ----------------------
% \subsection{See also}
% % ----------------------

% ----------------------
\subsection{1828 American Dictionary of the English Language} % *1828
% ----------------------
	\label{IDLENESS-1828}

	\begin{compactitem}
		\item[1] Not employed; unoccupied with business; inactive; doing nothing.
		\item[2] Slothful; given to rest and ease; averse to labor or employment; lazy; as an \textit{idle} man; an \textit{idle} fellow.
		\item[3] Affording leisure; vacant; not occupied; as \textit{idle} time; \textit{idle} hours.
		\item[4] Remaining unused; unemployed; applied to things; as, my sword or spear is \textit{idle}.
		\item[5] Useless; vain; ineffectual; as \textit{idle} rage.
		\item[6] Unfruitful; barren; not productive of good.
		\item[7] Trifling; vain; of no importance; as an \textit{idle} story; an \textit{idle} reason; \textit{idle} arguments.
		\item[8] Unprofitable; not tending to edification. Every \textit{idle} word that men shall speak, they shall give an account thereof in the day of judgment (\hyperref[MATTHEW-12-36]{Matthew 12:36}).
		\item Idle differs from lazy; the latter implying constitutional or habitual aversion or indisposition to labor or action, sluggishness; whereas \textit{idle} in its proper sense, denotes merely unemployed. An industrious man may be \textit{idle} but he cannot be lazy.
	\end{compactitem}

% % ----------------------
% \subsection{Oxford English Dictionary} % *OED
% % ----------------------
	% \label{IDLENESS-OED}

	% \begin{compactitem}
		% \item[]
	% \end{compactitem}

% ----------------------
\subsection{Scriptural Usage} % *SCRIPTURES
% ----------------------
	\label{IDLENESS-SCRIPTURES}

	% % -----
	% \subsubsection{Old Testament} % *OT
	% % -----
		% \label{IDLENESS-OT}
	
		% % --
		% \paragraph{KJV}
		% % --
			% \label{IDLENESS-OT-KJV}
	
			% \begin{tabular}{ll}
				% Total occurrences in the OT & 
			% \end{tabular}
			
			% \begin{compactdesc}
				% \item[]
			% \end{compactdesc}
			
		% % --
		% \paragraph{Hebrew Bible}
		% % --
			% \label{IDLENESS-HEBREW}
		
			% %
			% \subparagraph{\texorpdfstring{\he{sample word}}{}} % paste actual hebrew text here
			% %
				% \label{PAGA} % whatever is in step bible without periods
			
				% \begin{compactdesc}
					% \item[HALOT , PDF ] \textbf{} % Use the 1994 PDF
						% \begin{compactitem}
							% \item[1]
						% \end{compactitem}
					% \item[BDB , PDF ] \textbf{}
						% \begin{compactitem}
							% \item[1]
						% \end{compactitem}
					% \item[DCH PDF ] \textbf{} % don't include the actual page number, they repeat from volume to volume
						% \begin{compactitem}
							% \item[1]
						% \end{compactitem}
				% \end{compactdesc}
				
				% \begin{compactdesc}
					% \item[Some OT uses] \textbf{}
						% \begin{compactdesc}
							% \item[]
						% \end{compactdesc}
				% \end{compactdesc}
			
		% % --
		% \paragraph{Septuagint}
		% % --
			% \label{IDLENESS-LXX}
		
			% %
			% \subparagraph{\texorpdfstring{\gr{sample word}}{}}
			% %
		
	% -----
	\subsubsection{New Testament} % *NT
	% -----
		\label{IDLENESS-NT}
	
		% --
		\paragraph{KJV}
		% --
			\label{IDLENESS-NT-KJV}
	
			\begin{tabular}{ll}
				Total occurrences in the NT & 7
			\end{tabular}
			
			All the occurrences of English ``idle'' in the NT are \hyperref[ARGOS]{\grl{ἀργός}} except for Luke 24:11, which is \gr{λῆρος}, a noun, meaning ``nonsense, idle talk.''
			
			\begin{compactdesc}
				\item[Matthew 12:36] But I say unto you, That every idle word that men shall speak, they shall give account thereof in the day of judgment.
				\item[Matthew 20:3] And he went out about the third hour, and saw others standing idle in the marketplace,
				\item[Matthew 20:6] \textbf{(2)} And about the eleventh hour he went out, and found others standing idle, and saith unto them, Why stand ye here all the day idle?
				\item[Luke 24:11] And their words seemed to them as idle tales, and they believed them not.
					\begin{compactitem}
						\item ``Idle tales'' here is \gr{λῆρος}, ``nonsense, idle talk.''
					\end{compactitem}
				\item[1 Timothy 5:13] \textbf{(2)} And withal they learn to be idle, wandering about from house to house; and not only idle, but tattlers also and busybodies, speaking things which they ought not.
			\end{compactdesc}
			
		% --
		\paragraph{Greek New Testament} % *GREEK
		% --
			\label{IDLENESS-GREEK}
		
			%
			\subparagraph{\texorpdfstring{\gr{ἀργός}}{}}
			%
				\label{ARGOS}
			
				\begin{compactdesc}
					\item[BDAG 112a] \textbf{}
						\begin{compactitem}
							\item[1] \textbf{pertaining to being without anything to do, \textit{unemployed, idle}}
								\begin{compactitem}
									\item Matthew 20:3, 6 listed here, with 1 Timothy 5:13.
								\end{compactitem}
							\item[2] \textbf{pertaining to being unwilling to work, \textit{idle, lazy}}
								\begin{compactitem}
									\item The only NT use mentioned here is \hyperref[TITUS-1-12]{Titus 1:12}.
								\end{compactitem}
							\item[3] \textbf{pertaining to being unproductive, \textit{useless, worthless}}
								\begin{compactitem}
									\item James 2:20 is listed here, along with 2 Peter 1:8 and Matthew 12:36.
								\end{compactitem}
						\end{compactitem}
					\item[CGL 212b, PDF 238] \textbf{}
						\begin{compactitem}
							\item 
						\end{compactitem}
					\item[LSJ 236a, PDF 309] \textbf{}
						\begin{compactitem}
							\item 
						\end{compactitem}
				\end{compactdesc}
				
				\begin{tabular}{ll}
					Total occurrences in the NT & 8
				\end{tabular}
				
				\begin{compactdesc}
					\item[Some NT uses] \textbf{}
						\begin{compactdesc}
							\item[Matthew 12:36] \gr{λέγω δὲ ὑμῖν ὅτι πᾶν ῥῆμα ἀργὸν ὃ λαλήσουσιν οἱ ἄνθρωποι, ἀποδώσουσιν περὶ αὐτοῦ λόγον ἐν ἡμέρᾳ κρίσεως·}
							\item[Matthew 20:3] 
							\item[Mathtew 20:6] 
							\item[1 Timothy 5:13] \textbf{(2)}
							\item[Titus 1:12] \gr{εἶπέν τις ἐξ αὐτῶν, ἴδιος αὐτῶν προφήτης, Κρῆτες ἀεὶ ψεῦσται, κακὰ θηρία, γαστέρες ἀργαί·}
							\item[James 2:20] \gr{θέλεις θέλεις δὲ γνῶναι, ὦ ἄνθρωπε κενέ, ὅτι ἡ πίστις χωρὶς τῶν ἔργων ἀργή ἐστιν;}
								\begin{compactitem}
									\item The KJV translates \gr{νεκρά} here instead of \gr{ἀργή}; see \hyperref[JAMES-2-20-dead]{this note}.
								\end{compactitem}
							\item[2 Peter 1:8] \gr{ταῦτα γὰρ ὑμῖν ὑπάρχοντα καὶ πλεονάζοντα οὐκ ἀργοὺς οὐδὲ ἀκάρπους καθίστησιν εἰς τὴν τοῦ κυρίου ἡμῶν Ἰησοῦ Χριστοῦ ἐπίγνωσιν·}
						\end{compactdesc}
				\end{compactdesc}
		
	% % -----
	% \subsubsection{Book of Mormon} % *BOM
	% % -----
		% \label{IDLENESS-BOM}
	
		% \begin{tabular}{ll}
			% Total occurrences in the BOM & 
		% \end{tabular}
		
		% \begin{compactdesc}
			% \item[]
		% \end{compactdesc}
		
	% % -----
	% \subsubsection{Doctrine and Covenants} % *DC
	% % -----
		% \label{IDLENESS-DC}
	
		% \begin{tabular}{ll}
			% Total occurrences in the DC & 
		% \end{tabular}
		
		% \begin{compactdesc}
			% \item[]
		% \end{compactdesc}
		
	% % -----
	% \subsubsection{Pearl of Great Price} % *PGP
	% % -----
		% \label{IDLENESS-PGP}
	
		% \begin{tabular}{ll}
			% Total occurrences in the PGP & 
		% \end{tabular}
		
		% \begin{compactdesc}
			% \item[]
		% \end{compactdesc}
		
% % ----------------------
% \subsection{Key Phrases} % *KEYPHRASES *PHRASES
% % ----------------------
	% \label{IDLENESS-PHRASES}

	% \begin{compactdesc}
		% \item[] \textbf{\#times} | list
			% % \begin{compactdesc}
				% % \item[Bible anchor verses] \phantom{.}
					% % \begin{compactdesc}
						% % \item[Romans 5:5] \gr{}
					% % \end{compactdesc}
				% % \item[Commentary] ---
			% % \end{compactdesc}
	% \end{compactdesc}
	
% % ----------------------
% \subsection{Bible Dictionaries} % *DICTIONARIES
% % ----------------------
	% \label{IDLENESS-BD}

% % ----------------------
% \subsection{Restoration Usage} % *RESTORATION
% % ----------------------
	% \label{IDLENESS-RESTORATION}

	% % -----
	% \subsubsection{Joseph Smith} % *JS
	% % -----
		% \label{IDLENESS-JS}
	
		% \begin{compactdesc}
			% \item[{\hyperref[KEY]{Date/Source}}]
		% \end{compactdesc}
		
	% % -----
	% \subsubsection{Temple Ordinances}
	% % -----
		% \label{IDLENESS-TEMPLE}
	
		% \begin{compactdesc}
			% \item[{\hyperref[KEY]{Date/Source}}]
		% \end{compactdesc}
	
	% % -----
	% \subsubsection{Brigham Young} % *BY
	% % -----
		% \label{IDLENESS-BY}
	
		% \begin{compactdesc}
			% \item[{\hyperref[KEY]{Date/Source}}]
		% \end{compactdesc}
	
% % ----------------------
% \subsection{Commentary and Studies} % *COMMENTARY *STUDIES
% % ----------------------
	% \label{IDLENESS-COMMENTARY}

	% % -----
	% \subsubsection{Key Points}
	% % -----
		% \label{IDLENESS-KEYS}
	
		% \begin{compactitem}
			% \item
		% \end{compactitem}
		
	% % -----
	% \subsubsection{Sample Study}
	% % -----
		% \label{IDLENESS-study}